\section{The scoped one-parameter wrapper}
\label{sec:wrapper}

The one-parameter open domain is
\[
0<\theta\le120^\circ.
\]
The certificate uses the half-angle coordinate
\[
t=\tan(\theta/2),\qquad 0<t\le\sqrt3,
\]
and many sign proofs are carried out on the rational open superset
\[
0<t<2.
\]

Every representative tree has six unordered pair obligations.  The A7d wrapper
routes the twelve obligations across the two trees as follows.
\begin{center}
\begin{tabular}{lll}
\toprule
Tree & Pair & Route\\
\midrule
\treeA{} & \(P_0-P_1\) & B04 selected-hinge contact-side\\
\treeA{} & \(P_0-P_2\) & B04 selected-hinge contact-side\\
\treeA{} & \(P_0-P_3\) & B05 common-edge projection\\
\treeA{} & \(P_1-P_2\) & B05 common-edge projection\\
\treeA{} & \(P_1-P_3\) & B04 selected-hinge contact-side\\
\treeA{} & \(P_2-P_3\) & B06/B07 shared-face residual\\
\midrule
\treeB{} & \(P_0-P_1\) & B04 selected-hinge contact-side\\
\treeB{} & \(P_0-P_2\) & B06/B07 shared-face residual\\
\treeB{} & \(P_0-P_3\) & B05 common-edge projection\\
\treeB{} & \(P_1-P_2\) & B05 common-edge projection\\
\treeB{} & \(P_1-P_3\) & B04 selected-hinge contact-side\\
\treeB{} & \(P_2-P_3\) & B04 selected-hinge contact-side\\
\bottomrule
\end{tabular}
\end{center}

\begin{theorem}[Scoped zero-thickness one-parameter wrapper]
\label{thm:wrapper}
For the catalogued median-plane S4 pieces, the public package records a closed
zero-thickness one-parameter wrapper for the two audited representatives
\treeA{} and \treeB{} on the domain
\[
\theta=0\quad\text{and}\quad 0<\theta\le120^\circ.
\]
At \(\theta=0\), the configuration is the closed-contact endpoint of
Lemma~\ref{lem:endpoint}.  On the open ray, every unordered pair obligation in
each representative tree is assigned to a completed certificate layer: six B04
selected-hinge bridge rows, four B05 common-edge projection rows, and two
B06/B07 shared-face residual rows.  There are no route-clean B03 obligations on
this one-parameter spine.
\end{theorem}

\begin{proof}
The endpoint clause follows from Lemmas~\ref{lem:tiling}--\ref{lem:kinematics}:
at \(\theta=0\), all rooted transforms are identity and the pieces are in the
catalogued closed tiling.

For the open ray, the A7d manifest is a pair-route ledger.  It contains two
wrapper records, one for \treeA{} and one for \treeB{}, each with six unordered
piece-pair obligations.  The manifest records the predicate route counts
\[
\text{B04}=6,
\qquad
\text{B05}=4,
\qquad
\text{B06/B07}=2.
\]
The B03 count is zero by the A7b route-vacuity certificate.  The routed layers
are discharged by the certified propositions in Sections~\ref{sec:layers} and
\ref{sec:b04}.  Thus every open-ray pair obligation is routed to a completed
package-local layer, and no additional B03 obligation remains on this scoped
one-parameter domain.
\end{proof}

\begin{certprop}[A7d manifest]
The theorem wrapper above is the statement recorded in
\path{certified/a7d_one_parameter_theorem_wrapper_manifest.json}.  The manifest
has \texttt{record\_count=2}, \texttt{wrapper\_closed\_count=2}, and object
status count
\texttt{a7d\_one\_parameter\_ray\_wrapper\_closed: 2}.
\end{certprop}

\begin{remark}[Scope]
The theorem is a route-wrapper theorem for a zero-thickness one-parameter model.
It is not a positive-thickness fabrication theorem, not a global collision-free
motion theorem, and not a three-parameter bounded-cell theorem.
\end{remark}
